\section{Related Work}

\subsection{Execution Feedback for Code Generation}

Reinforcement learning from execution feedback has emerged as a dominant paradigm for aligning code generation models. PPOCoder \citep{shojaee2023ppo} applies Proximal Policy Optimization with binary test pass/fail rewards, demonstrating that execution feedback can improve code generation without human annotation. CodeRL \citep{le2022coderl} extends this with actor-critic methods and self-repair mechanisms. InterCode \citep{yang2023intercode} formalizes interactive coding as an RL environment where code serves as actions and execution feedback as observations.

Recent work has explored richer reward signals. RLPF \citep{jing2026rlpf} introduces staged rewards that order failed programs by execution progress and rank correct programs by efficiency improvement. ReTool \citep{feng2025retool} integrates real-time code execution within reasoning, using outcome feedback to guide tool use.

Our work is orthogonal to reward design: we focus on \emph{how} credit is assigned to individual tokens given any reward signal, rather than \emph{what} reward signal is used.

\subsection{Fine-Grained Credit Assignment}

StepCoder \citep{dou2024stepcoder} introduces Fine-Grained Optimization (FGO), which masks non-executed code segments from gradient updates. However, StepCoder introduces FGO alongside Curriculum of Code Completion Subtasks (CCCS), making it impossible to isolate which component drives improvement.

Fine-Grained RLHF \citep{wu2023finegrained} demonstrates that per-segment rewards outperform holistic rewards in general language tasks, supporting the value of dense credit assignment.

Our work provides the first controlled validation of FGO's mechanism, isolating it from curriculum effects.

\subsection{Inference-Time Execution Guidance}

EG-CFG \citep{lavon2025egcfg} achieves state-of-the-art results (99.4\% on HumanEval) by injecting runtime feedback directly into the decoding process. While impressive, inference-time methods do not improve the underlying model's capabilities.

Our focus is training-time mechanism validation.
